% Standalone source-record documentation; no research attempts.
\documentclass[11pt]{article}
\usepackage[margin=25mm]{geometry}
\usepackage{fontspec,unicode-math}
\setmainfont{Latin Modern Roman}
\setmathfont{Latin Modern Math}
\usepackage{amsmath,amssymb,mathtools,xurl,hyperref}
\hypersetup{colorlinks=true,urlcolor=blue,linkcolor=blue}
\setcounter{secnumdepth}{0}
\setlength{\parindent}{0pt}
\setlength{\parskip}{5pt}
\setlength{\emergencystretch}{3em}
\newcommand{\sref}[2]{\hyperlink{source:#1:#2}{[S#2]}}
\newcommand{\cataloguescope}[1]{\paragraph{Source record.}#1}
\newcommand{\floor}[1]{\left\lfloor#1\right\rfloor}
\begin{document}
% UnsolvedMath ID 20001172; source catalogue code AIM-COMPUTATION-0010
\setcounter{section}{20001171}
\section{Certified weak-coupling tubes for structured reachability}\label{20001172}
{\small\sffamily UnsolvedMath ID 20001172 · Computer Science}
\subsection{Definitions and mathematical statement}

\paragraph{Shared notation.}Standard number systems are $\mathbb N,\mathbb Z,\mathbb Q,\mathbb R,\mathbb C$. Specialized notation and index ranges follow the cited source and the definitions below.


For dynamics $\dot x=f(x,u,w)$, controls $u$ and disturbances $w$ determine reachable states. Inner and outer approximations respectively certify membership and contain all possible reachable states. A level-set representation describes a set by a sublevel set of a scalar function. Weak coupling means interactions between components are small in a stated quantitative model; probabilistic verification requires a sampling law and a confidence guarantee.
\paragraph{Record scope and qualifications.}
The catalogue title is retained as metadata. Named certificates, contracts, bounds or frameworks in that title are not asserted as results of the original source extract.
This is a family of numerical-verification questions. The dynamics, coupling scale, horizon, safe set, sampling law, and sense of certification are not fixed; no specific tube theorem is inferred. The supplied catalogue marks this record open; its present research status was not independently checked.
The source text below is rendered literally where necessary; damaged formulas and composite extracts are not silently repaired into a new mathematical claim.
\paragraph{Statement recorded in the supplied source.}
A.7 Mitchell, Ian

The solution to the Hamilton-Jacobi(-Bellman)(-Isaacs) partial differential equation (HJ PDE) provides the optimal value function (eg cost to go) for bounded input optimal control or zero-sum differential game problems respectively. These equations arise in many disparate fields, and my personal interest lies in how they can be used to approximate the reachable sets of continuous dynamic systems in order to verify and/or synthesize provably safe control policies. For these purposes, it may be possible to relax the goal of a convergent approximation of the solution of the PDE. Instead, we seek guaranteed under or over approximations of the solutions; for example, an underapproximation of the solution might overapproximate the unsafe reachable set, which might thereby generate false negative verifications but no false positives. With or without such relaxations, is it possible to take advantage of structure in the system dynamics to reduce the typically exponential growth in computational cost with system dimension? For example, is it possible to take advantage of weak or linear coupling between certain components to reduce the degrees of freedom needed by an approximation, perhaps through partial parameterization? Or another approach: Can sampling methods be applied to produce quantifiable probabilistic verifications? In addition to these questions related to the application area of verification, I am inter-ested in the class of algorithms designed to approximate solutions of HJ PDEs - often broadly called "level set methods." This interest includes both the time-dependent and stationary (or static) versions of the equation, as well as issues related to the use of such numerical approximations in the control of robotic systems. \sref{20001172}{2}
\paragraph{Source status.}
The supplied catalogue labels this source record open; that metadata was checked against the supplied snapshot on 2026-10-11. The label is preserved for provenance and is not an independent verification of present open status or mathematical correctness.
\subsection{Short English statement}
The source asks whether structure or sampling can reduce the cost of certified reachability computations.
\subsection{Sources}
\begin{itemize}
\item[\hypertarget{source:20001172:1}{S1}] Ulam AI and the UnsolvedMath Contributors, UnsolvedMath: A Curated Collection of Open Mathematics Problems, record 20001172 (AIM-COMPUTATION-0010). Catalogue curation and metadata: CC BY 4.0. The source extract is retained; definitions, English summaries and qualifications are editorial additions.
\url{https://huggingface.co/datasets/ulamai/UnsolvedMath}.
\item[\hypertarget{source:20001172:2}{S2}] AIM Workshop Problem Lists, source linked by catalogue record AIM-COMPUTATION-0010.
\url{https://aimath.org/WWN/optimalcontrol/optimalcontrol.pdf}.
\end{itemize}
\paragraph{Preparation and review.}
The original source and UnsolvedMath contributors retain attribution for the source extract. Codex (GPT-6; exact serving revision not exposed) assisted with definitions, English summaries, scope notes and formatting. No new human mathematical authorship or independently verified solution is asserted. Permanent source ID 20001172 and code AIM-COMPUTATION-0010 checked against the complete supplied registry and current repository. Distinct upstream records, including repeated questions and broad or damaged extracts, are retained for complete coverage of the source CS category. A repeated record is not a new independent conjecture. Complete mathematical-equivalence and current-literature review has not been performed. Preparation and final review follow contribution prompts C and D, with source gaps disclosed. Structural and rendering checks do not establish a new theorem or open status.
\label{end:20001172}
\end{document}
